\subsection{Motivation}

Despite all of its merits, achieving the best antenna orientation in practice still poses significant challenges. The classical approaches to automatic alignment of the antenna involve exhaustive scanning and hill-climbing using the received signal strength (RSSI). They are widely used due to their simplicity. While scanning all possible orientations of the antenna takes too much time and energy, hill climbing algorithms are prone to converge at local maxima because of the noisy RSSI.

Recent developments have focused on employing machine learning and reinforcement learning for antenna alignment and beam selection mainly in applications dealing with mmWave and phased array antennas ~\cite{ref1,ref2}. Despite high efficiency, all of those approaches require sophisticated radio frequency hardware, vast amounts of data and large computations. Therefore, they cannot be applied to low cost, sub-6 GHz systems due to the lack of necessary resources. Among various methods of reinforcement learning, Q-learning appears especially suitable for constrained systems since it does not require any prior knowledge of the environment and needs no large datasets for efficient operation.

In this project, we aim to evaluate the feasibility of using simple \textbf{Q-learning} algorithms along with RSSI and mechanical steering for improving antenna alignment as compared to traditional algorithms.